\begin{figure}[!ht]
  \centering
  \includegraphics[width=\linewidth]{figures/teaser_figure.pdf}
  \caption{
    \textbf{\methodname.} Given multiple input views \textbf{(top-left)}, \methodname{} reconstructs camera poses and consistent depth by repeatedly applying the \emph{same} transformer block, with the number of refinement steps $K$ exposed as an inference-time compute knob. Decoding the intermediate state of a single $K{=}16$ forward pass at iterations $k \in \{2, 4, 8, 16\}$ shows progressively sharper geometry and more accurate camera poses (\textbf{right}; frustums are colored by per-camera error). Across five benchmarks \textbf{(bottom-left)}, \methodname{} matches or surpasses much larger feed-forward baselines at a small fraction of their parameter count (dot area).
  }
  \label{fig:teaser}
\end{figure}


\begin{abstract}
Recent feed-forward 3D reconstruction transformers have scaled to over a billion parameters, following the broader trend of increasing model capacity in computer vision. Yet emerging evidence suggests that contiguous transformer layers often behave like repeated applications of similar operations~\citep{jacobs2025raptor}, and multi-view reconstruction transformers refine their predictions progressively across decoder depth~\citep{stary2025understanding}. We posit that model depth partially buys iteration, paid for inefficiently in unique parameters, and instead make that iteration explicit in architecture. Our model, \methodname, applies a single looped transformer block recurrently to per-view features for $K$ refinement steps. Trained once, it exposes $K$ as an inference-time compute knob, matching or outperforming substantially larger feed-forward baselines across five reconstruction benchmarks spanning indoor, outdoor, object-centric, and driving scenes, while using a fraction of their parameters and comparable or lower compute. Importantly, the same looped block formulation outperforms an otherwise identical variant with independent per-step parameters under matched training data and compute, suggesting that explicit iteration is not merely a compute-efficient substitute for capacity but a stronger inductive bias for multi-view 3D reconstruction. We release code and model on our \href{https://research.nvidia.com/labs/dvl/projects/dvlt/}{project page https://research.nvidia.com/labs/dvl/projects/dvlt/}.
\end{abstract}
\vspace{-6mm}
